\documentclass[11pt]{article}
\usepackage{shared/hanzocover}
\usepackage[utf8]{inputenc}
\usepackage[T1]{fontenc}
\usepackage{amsmath,amssymb}
\usepackage{graphicx}
\usepackage{booktabs}
\usepackage{xcolor}
\usepackage[hidelinks]{hyperref}
\emergencystretch=1.5em

\title{Context Engine:\\ Typed Retrieval Across Memory and Code}
\author{Zach Kelling\thanks{Corresponding author: research@hanzo.ai} \\
    \textit{Hanzo Industries Inc (Techstars '17)} \\
    Los Angeles, California \\
    \texttt{https://hanzo.ai}}
\date{Draft, revision~0}

\begin{document}

\hanzocoverpage

\twocolumn[
  \begin{@twocolumnfalse}
  \maketitle
  \begin{abstract}
  \noindent
  The companion paper showed that long-term conversational memory is a
  linkage problem: a question reaches its evidence through a fact, an
  entity or a date more often than through wording, and a second index of
  typed links lifts complete-evidence recall on multi-hop questions from
  22.7\% to 36.2\% with nothing learned. This paper asks whether the same
  engine, unchanged in shape, is the right retriever for a software
  repository. A repository is connected information with an unusually
  honest ground truth: the compiler and the language server know which
  symbol a reference resolves to, which function calls which, which test
  covers which file, which commit last touched a line. We describe the
  Context Engine as one abstraction---raw artefact, typed facts, canonical
  entities, time, provenance, and a read path that plans, seeds, expands
  along typed edges within a budget, scores once and rewrites the query
  after each resolved hop---and its two instantiations: conversation, where
  the edges are people, events, sessions and dates; and code, where they are
  symbols, definitions, references, implementations, callers, callees,
  types, files, commits, issues, tests and runtime logs. We set out a
  deterministic benchmark whose gold links are emitted by the language
  server rather than annotated by hand, and an evaluation plan on RepoBench,
  CrossCodeEval and RepoProbe with retrieval measured separately from
  generation. Results are pending and the sections that will carry them are
  marked. The claim this draft is written to test is narrow: that semantic
  similarity is one candidate generator inside a typed retrieval system,
  and that the system beats similarity alone on both kinds of corpus.
  \end{abstract}
  \vspace{1em}
  \end{@twocolumnfalse}
]

\section{One Abstraction}

Retrieval systems for memory and retrieval systems for code are usually
built by different teams from different papers. Both start from the same
mistake: embed the unit of text, embed the request, return the nearest
neighbours, and hope the answer resembles the question. In a conversation
the turn that answers ``where did she work after leaving Acme'' contains
neither ``work'' nor ``Acme''. In a repository the file that explains why a
request handler retries contains neither ``retry'' nor the handler's name;
it is the interface the handler implements, two directories away, reached
by a typed edge the compiler already knows.

The Context Engine is the abstraction the companion paper
\cite{contextmemory} arrived at for memory, restated so that nothing in it
names a conversation:

\begin{itemize}
  \item \textbf{raw} --- the artefact as received, immutable, with source and
        observation time.
  \item \textbf{facts[]} --- atomic typed statements about the artefact, each
        with a canonical subject, a relation, a canonical object or a literal,
        a time interval where one applies, a confidence, an embedding and a
        pointer to the exact span it came from.
  \item \textbf{entities[]} --- canonical identifiers resolved by type.
  \item \textbf{events[]} --- groups of facts that belong together.
  \item \textbf{supersedes[]} --- what this artefact replaces.
  \item \textbf{abstraction} and \textbf{anchors[]} --- one short phrase for what
        the artefact is about and several alternative entry points.
\end{itemize}

And the read path: understand the request into operators; seed from dense,
lexical, fact, entity and time generators with small pools; expand along
typed edges within a bounded budget; score every candidate once with
recorded contributions; return the least raw context that proves the
answer; and if a hop resolved something, rebuild the query from the resolved
state before the next hop.

\section{Two Instantiations}

\begin{table*}[t]
\centering
\footnotesize
\begin{tabular}{p{24mm}p{58mm}p{74mm}}
\toprule
Record field & Conversation & Code \\
\midrule
raw & a turn & a file, a hunk, a log line \\
entity & person, place, org, work, date & symbol, type, module, file, commit, issue, test \\
relation & said, did, plans, moved, met & defines, references, implements, calls, imports, extends, covers, touched \\
event & the trip, the interview & the change set, the incident, the release \\
time & session date, event time & commit time, blame, run time \\
supersedes & a later fact about the same relation & a later commit on the same lines \\
abstraction & ``Caroline moved from Sweden'' & ``retry policy for outbound HTTP'' \\
anchors & names, places, dates & identifiers, error strings, config keys \\
provenance & conversation, session, turn & repo, path, line range, commit \\
\bottomrule
\end{tabular}
\caption{The same record, two corpora. Every edge in the code column is
emitted by tools that already exist: the language server, the compiler,
the version control system, the test runner.}
\label{tab:two}
\end{table*}

\subsection{Conversation}

The memory instantiation is described and measured in the companion paper.
Its edges are extracted by a model and are probabilistic; its gold evidence
is human-annotated and, as LoCoMo-Conv showed \cite{locomoconv},
incomplete. This is the harder of the two settings to evaluate honestly,
which is why the companion paper spends a section on what is being measured.

\subsection{Code}

The code instantiation is the easier of the two to evaluate and the harder
to build, because the typed edges are many and exact. Candidate generators:

\begin{itemize}
  \item dense search over facts about symbols (a symbol's abstraction: its
        signature, docstring and the one-line purpose an extractor assigns);
  \item lexical search over identifiers and literal strings, which in code
        carries more of the signal than in prose;
  \item the symbol index: exact matches on canonical symbol identifiers, with
        aliases and re-exports resolved;
  \item typed edges from the language server: definition, references,
        implementations, type hierarchy, callers and callees;
  \item file and module structure: siblings, parents, the AST path;
  \item history: the commits that touched a span, their messages, the issues
        they close, blame;
  \item tests: the test cases whose coverage includes a span, and the runtime
        logs a failing test emitted.
\end{itemize}

Expansion is bounded exactly as in memory---at most $n$ hops along at most
$m$ edges of each type---and every candidate carries the edge it entered by.
A stateful hop is natural here: resolving ``the retry logic'' to a symbol on
hop one turns hop two into \emph{implementations of that symbol's
interface, in files touched by the commit that introduced it}, a query that
no embedding of the original sentence would find.

\section{A Deterministic Gold Benchmark}
\label{sec:lsp}

Hand-annotated retrieval gold is expensive and incomplete. For code it is
also unnecessary. We build the Hanzo LSP benchmark as follows. For each of
a set of repositories, run the language server and the test runner and
record every definition, reference, implementation, call, import and
coverage edge. Generate questions mechanically from the graph---``what
implements X'', ``who calls Y'', ``which test covers Z'', ``what changed in
this function in the last release''---and their answers as the exact set of
spans the graph names. Phrase each question three ways, one of them without
the symbol's name, so that lexical matching alone cannot answer it. Score
retrieval as ALL and ANY over the gold spans at $k$, as in the companion
paper, and report MRR, tokens delivered and latency.

The gold is complete by construction, so the SUPPORTED grade of the memory
paper is not needed here; EXACT is the whole story. This is the benchmark on
which a difference between similarity and linkage will be least arguable.

\section{External Lanes}

\paragraph{RepoBench.} \cite{repobench} separates retrieval (RepoBench-R),
completion (RepoBench-C) and the pipeline (RepoBench-P) for Python and Java.
We report RepoBench-R as a retrieval number and RepoBench-P with a fixed
generator across rows, so a gain is attributable to retrieval.

\paragraph{CrossCodeEval.} \cite{crosscodeeval} requires cross-file context
to complete code in Python, Java, TypeScript and C\#. We report exact match
and edit similarity with the same generator across rows and the retrieval
budget held constant.

\paragraph{RepoProbe.} \cite{repoprobe} asks open-ended architectural
questions drawn from GitHub Discussions across real repositories and scores
answers against checklists of atomic verifiable facts. We report the
checklist score and, separately, the fraction of checklist facts whose
supporting span was retrieved, which is the retrieval question this paper
is about.

\paragraph{ContextBench.} \cite{contextbench} measures context retrieval in
coding agents directly and shares a name with the Hanzo benchmark described
in the companion paper; the two are distinct, and we run this one as an
external lane where its harness permits.

\section{Evaluation Plan}

Every table carries the rows of the companion paper's protocol: semantic
only; plus lexical; plus facts; plus entities; plus timeline; plus adjacency
(here: siblings in the file and the AST); plus typed graph; plus iterative
hops; full; and the failed variants---pseudo-relevance feedback, generic
multi-query, chain search, surface-identifier expansion, global fusion.
Retrieval is measured separately from generation. A declared dev set of
repositories is used for every parameter choice; the commit is frozen; the
test repositories are run once.

\section{Results}
\label{sec:results}

Pending. The tables below will be filled from \texttt{runs/} of the benchmark
repository with the run identifier printed beside each.

\begin{table*}[t]
\centering
\small
\begin{tabular}{lrrrr}
\toprule
Hanzo LSP benchmark & ALL@10 & ANY@10 & MRR & tokens/q \\
\midrule
semantic only & \ldots & \ldots & \ldots & \ldots \\ % TODO run code-lsp-test-semantic
+ lexical & \ldots & \ldots & \ldots & \ldots \\
+ symbol index & \ldots & \ldots & \ldots & \ldots \\
+ typed edges & \ldots & \ldots & \ldots & \ldots \\
+ history and tests & \ldots & \ldots & \ldots & \ldots \\
+ iterative hops & \ldots & \ldots & \ldots & \ldots \\
full & \ldots & \ldots & \ldots & \ldots \\
\bottomrule
\end{tabular}
\caption{Pending: retrieval on the deterministic benchmark of
Section~\ref{sec:lsp}.}
\end{table*}

\begin{table*}[t]
\centering
\small
\begin{tabular}{lrrr}
\toprule
External lane & retrieval & end-to-end & tokens/q \\
\midrule
RepoBench-R / P & \ldots & \ldots & \ldots \\ % TODO run code-repobench-test-full
CrossCodeEval & \ldots & \ldots & \ldots \\ % TODO run code-cceval-test-full
RepoProbe & \ldots & \ldots & \ldots \\ % TODO run code-repoprobe-test-full
\bottomrule
\end{tabular}
\caption{Pending: external lanes, full configuration, one generator across rows.}
\end{table*}

\section{Related Work}

Repository-level retrieval for completion is the subject of RepoBench
\cite{repobench} and CrossCodeEval \cite{crosscodeeval}; architectural
comprehension of RepoProbe \cite{repoprobe}; context retrieval for coding
agents of ContextBench \cite{contextbench}. On the memory side the closest
designs are AnchorMem \cite{anchormem}, which anchors retrieval on atomic
facts and returns raw context, Memora \cite{memora}, which indexes an
abstraction and cue anchors and keeps the value unindexed, and APEX-MEM
\cite{apexmem}, an entity-centric temporal property graph with append-only
history. HippoRAG \cite{hipporag} is the nearest graph-retrieval ancestor.
None of these is applied to code, and none of the code benchmarks is
approached with a memory system; the point of this draft is that the gap is
an accident of who built what.

\section{What Would Falsify It}

If the full configuration does not beat semantic-only retrieval on the
deterministic benchmark by a margin outside the bootstrap interval, the
abstraction is wrong for code and the paper says so. If it wins there but
not on RepoBench-P with a fixed generator, the gain is not reaching the
task. If it wins on code but the memory numbers of the companion paper do
not hold on their test split, the engine is two engines with one name.
Each of these is a reportable outcome.

\begin{thebibliography}{99}
\bibitem{contextmemory} Z. Kelling. Linkage, Not Similarity: Context-Sensitive Retrieval
  for Long-Term Agent Memory. Hanzo Research, 2026.
\bibitem{locomoconv} W.-Y. Chang, Y.-N. Chen. When Users Don't Ask: Benchmarking
  Context-Driven Memory Retrieval in Conversational Agents. arXiv:2609.03467, September 2026.
\bibitem{repobench} T. Liu, C. Xu, J. McAuley. RepoBench: Benchmarking Repository-Level
  Code Auto-Completion Systems. arXiv:2306.03091, 2023.
\bibitem{crosscodeeval} Y. Ding, Z. Wang, W. U. Ahmad, H. Ding, M. Tan, N. Jain,
  M. K. Ramanathan, R. Nallapati, et al. CrossCodeEval: A Diverse and Multilingual
  Benchmark for Cross-File Code Completion. arXiv:2310.11248, 2023.
\bibitem{repoprobe} Y. Yang, A. Wu, J. Luo, R. Xuan, Z. Hu, Y. Liu, Z. Qin. RepoProbe:
  Benchmarking Architecture-Aware Repository Comprehension with Checklists.
  arXiv:2608.04783, August 2026; ASE 2026.
\bibitem{contextbench} H. Li, L. Zhu, B. Zhang, R. Feng, J. Wang, Y. Pan, E. T. Barr,
  F. Sarro, et al. ContextBench: A Benchmark for Context Retrieval in Coding Agents.
  arXiv:2602.05892, 2026.
\bibitem{anchormem} Z. Shen, S. Cheng, Z. Guo, W. Wang, Y. Wang, H. Huang. AnchorMem:
  Anchored Facts with Associative Contexts for Building Memory in Large Language Models.
  arXiv:2604.17377, April 2026.
\bibitem{memora} M. Xia, X. Zhang, S. Dixit, P. Harimurugan, R. Wang, V. R\"uhle, R. Sim,
  C. Bansal, et al. Memora: A Harmonic Memory Representation Balancing Abstraction and
  Specificity. arXiv:2602.03315, 2026; ICML 2026.
\bibitem{apexmem} P. Banerjee, M. Moshtaghi, S. Subramanian, A. Misra, A. Chadha. APEX-MEM:
  Agentic Semi-Structured Memory with Temporal Reasoning for Long-Term Conversational AI.
  arXiv:2604.14362, April 2026; ACL 2026.
\bibitem{hipporag} B. Jim\'enez Guti\'errez, Y. Shu, Y. Gu, M. Yasunaga, Y. Su. HippoRAG:
  Neurobiologically Inspired Long-Term Memory for Large Language Models. arXiv:2405.14831, 2024.
\end{thebibliography}

\end{document}
